\documentclass[10pt,a4paper]{amsart}
\usepackage{amsmath,amssymb,amsthm,graphicx}
% Retained preamble for the Stacks Project chapter snapshot in this directory.
% The chapter's own preamble.tex is a copy of the shared Stacks Project preamble; the block
% below is sliced MECHANICALLY, line for line, from it. Nothing was added, removed or
% redefined; the wrapper class, the xr-hyper/hyperref wiring and the externaldocument list
% are omitted because the assembler supplies the document class.
%
% lines taken from preamble.tex: 3, 5, 9, 11, 12, 156, 157, 158, ... (51 lines total)
%
% package coverage restored from the chapter's own preamble (2026-10-13): xy, xy-2cell

\usepackage[all]{xy}
\xyoption{2cell}
\UseAllTwocells

\newenvironment{reference}{\comment}{\endcomment}
\newenvironment{slogan}{\comment}{\endcomment}
\newenvironment{history}{\comment}{\endcomment}
\theoremstyle{plain}
\newtheorem{theorem}[subsection]{Theorem}
\newtheorem{proposition}[subsection]{Proposition}
\newtheorem{lemma}[subsection]{Lemma}
\theoremstyle{definition}
\newtheorem{definition}[subsection]{Definition}
\newtheorem{example}[subsection]{Example}
\newtheorem{exercise}[subsection]{Exercise}
\newtheorem{situation}[subsection]{Situation}
\theoremstyle{remark}
\newtheorem{remark}[subsection]{Remark}
\newtheorem{remarks}[subsection]{Remarks}
\numberwithin{equation}{subsection}
\def\lim{\mathop{\mathrm{lim}}\nolimits}
\def\colim{\mathop{\mathrm{colim}}\nolimits}
\def\Spec{\mathop{\mathrm{Spec}}}
\def\Hom{\mathop{\mathrm{Hom}}\nolimits}
\def\Ext{\mathop{\mathrm{Ext}}\nolimits}
\def\SheafHom{\mathop{\mathcal{H}\!\mathit{om}}\nolimits}
\def\SheafExt{\mathop{\mathcal{E}\!\mathit{xt}}\nolimits}
\def\Sch{\mathit{Sch}}
\def\Mor{\mathop{\mathrm{Mor}}\nolimits}
\def\Ob{\mathop{\mathrm{Ob}}\nolimits}
\def\Sh{\mathop{\mathit{Sh}}\nolimits}
\def\NL{\mathop{N\!L}\nolimits}
\def\CH{\mathop{\mathrm{CH}}\nolimits}
\def\proetale{{pro\text{-}\acute{e}tale}}
\def\etale{{\acute{e}tale}}
\def\QCoh{\mathit{QCoh}}
\def\Ker{\mathop{\mathrm{Ker}}}
\def\Im{\mathop{\mathrm{Im}}}
\def\Coker{\mathop{\mathrm{Coker}}}
\def\Coim{\mathop{\mathrm{Coim}}}
\DeclareMathSymbol{\boxtimes}{\mathbin}{AMSa}{"02}
\def\QCohstack{\mathcal{QC}\!\mathit{oh}}
\def\Cohstack{\mathcal{C}\!\mathit{oh}}
\def\Spacesstack{\mathcal{S}\!\mathit{paces}}
\def\Quotfunctor{\mathrm{Quot}}
\def\Hilbfunctor{\mathrm{Hilb}}
\def\Curvesstack{\mathcal{C}\!\mathit{urves}}
\def\Polarizedstack{\mathcal{P}\!\mathit{olarized}}
\def\Complexesstack{\mathcal{C}\!\mathit{omplexes}}
\def\Pic{\mathop{\mathrm{Pic}}\nolimits}
\def\Picardstack{\mathcal{P}\!\mathit{ic}}
\def\Picardfunctor{\mathrm{Pic}}
\def\Deformationcategory{\mathcal{D}\!\mathit{ef}}
\usepackage{cleveref}
\crefname{originallemma}{original lemma}{original lemmas}
\Crefname{originallemma}{Original lemma}{Original lemmas}
\makeatletter
\newlabel{lemma-characterize-quasi-coherent}{{lemma-characterize-quasi-coherent}{source}{}{sourcepointer-1}{}}
\newlabel{lemma-characterize-quasi-coherent@cref}{{[originallemma][0][]lemma-characterize-quasi-coherent}{[1][0][]source}{}{}{}}
\makeatother
\crefname{originallemma}{original lemma}{original lemmas}
\Crefname{originallemma}{Original lemma}{Original lemmas}
\makeatletter
\newlabel{lemma-pullback-quasi-coherent}{{lemma-pullback-quasi-coherent}{source}{}{sourcepointer-2}{}}
\newlabel{lemma-pullback-quasi-coherent@cref}{{[originallemma][0][]lemma-pullback-quasi-coherent}{[1][0][]source}{}{}{}}
\makeatother
\crefname{originallemma}{original lemma}{original lemmas}
\Crefname{originallemma}{Original lemma}{Original lemmas}
\makeatletter
\newlabel{lemma-stackification}{{lemma-stackification}{source}{}{sourcepointer-3}{}}
\newlabel{lemma-stackification@cref}{{[originallemma][0][]lemma-stackification}{[1][0][]source}{}{}{}}
\makeatother
\crefname{originallemma}{original lemma}{original lemmas}
\Crefname{originallemma}{Original lemma}{Original lemmas}
\makeatletter
\newlabel{sites-modules-lemma-local-final-object}{{03DN}{source}{}{sourcepointer-4}{}}
\newlabel{sites-modules-lemma-local-final-object@cref}{{[originallemma][0][]03DN}{[1][0][]source}{}{}{}}
\makeatother
\title{Stacks, Sheaves on Algebraic Stacks: lemma-stackification-quasi-coherent}
\author{The Stacks Project Authors}
\date{stacks-project:stacks-sheaves.tex:local-snapshot-20261012; GNU FDL 1.2 or later}
\begin{document}
\maketitle
\noindent\textit{Editorial note.} Statement span is lines 2196-2207 of sources/stacks-algebraic-stacks-sheaves/pinned-original-source.tex; proof span is lines 2209-2233, the complete author proof (23 body lines).\par\medskip
\noindent\textit{Editorial note.} Navigation stubs: 4 cited-but-not-included labels are supplied as non-load-bearing pointers with their pinned tags where the snapshot tag index has them.\par\medskip
\noindent\textit{Editorial note.} extra\_packages is empty on purpose: the retained preamble already loads xy with options and a second declaration clashes in the assembler wrapper.\par\medskip
\noindent\textit{Editorial note.} Text is retained verbatim, including the author's own typos and placeholders.\par\medskip
\noindent\textit{Editorial note.} Queue index 22 of queue\_algebraic-stacks-sheaves.json (pb=1210, stmt 2196-2207, proof 2209-2233).\par\medskip
\noindent\textit{Allowed standard prerequisites (editorial).}\begin{itemize}
\item Explicit prerequisite boundary. The cited inputs are prerequisites; the argument in the retained span is not replaced by any cited result.
\end{itemize}\medskip
\ifdefined\hypertarget\hypertarget{sourcepointer-1}{}\fi
\noindent\textit{Source reference.} Original lemma lemma-characterize-quasi-coherent: sources/stacks-algebraic-stacks-sheaves/pinned-original-source.tex, source lines 1798--1805. Cited by the retained statement or proof but not part of the collected spans; supplied only as a navigation pointer. The argument does not depend on its content.\par\medskip
\ifdefined\hypertarget\hypertarget{sourcepointer-2}{}\fi
\noindent\textit{Source reference.} Original lemma lemma-pullback-quasi-coherent: sources/stacks-algebraic-stacks-sheaves/pinned-original-source.tex, source lines 1777--1784. Cited by the retained statement or proof but not part of the collected spans; supplied only as a navigation pointer. The argument does not depend on its content.\par\medskip
\ifdefined\hypertarget\hypertarget{sourcepointer-3}{}\fi
\noindent\textit{Source reference.} Original lemma lemma-stackification: sources/stacks-algebraic-stacks-sheaves/pinned-original-source.tex, source lines 2173--2180. Cited by the retained statement or proof but not part of the collected spans; supplied only as a navigation pointer. The argument does not depend on its content.\par\medskip
\ifdefined\hypertarget\hypertarget{sourcepointer-4}{}\fi
\noindent\textit{Source reference.} Original lemma 03DN: sources/stacks-algebraic-stacks-sheaves/pinned-original-source.tex, source lines 2233--2233. Cited by the retained statement or proof but not part of the collected spans; supplied only as a navigation pointer. The argument does not depend on its content.\par\medskip
\begin{lemma}
\label{lemma-stackification-quasi-coherent}
Let $f : \mathcal{X} \to \mathcal{Y}$ be a $1$-morphism of categories
fibred in groupoids over $(\Sch/S)_{fppf}$. If
$f$ induces an equivalence of stackifications, then $f^*$
induces equivalences
$\textit{Mod}(\mathcal{O}_\mathcal{X}) \to
\textit{Mod}(\mathcal{O}_\mathcal{Y})$
and
$\QCoh(\mathcal{O}_\mathcal{X}) \to
\QCoh(\mathcal{O}_\mathcal{Y})$.
\end{lemma}

\begin{proof}
We may assume $\mathcal{Y}$ is the stackification of $\mathcal{X}$.
The first assertion is clear from
Lemma \ref{lemma-stackification}
and
$\mathcal{O}_\mathcal{X} = f^{-1}\mathcal{O}_\mathcal{Y}$.
Pullback of quasi-coherent sheaves are quasi-coherent, see
Lemma \ref{lemma-pullback-quasi-coherent}.
Hence it suffices to show that if $f^*\mathcal{G}$ is
quasi-coherent, then $\mathcal{G}$ is.
To see this, let $y$ be an object of $\mathcal{Y}$.
Translating the condition that $\mathcal{Y}$ is the stackification
of $\mathcal{X}$ we see there exists an fppf covering $\{y_i \to y\}$
in $\mathcal{Y}$ such that $y_i \cong f(x_i)$ for some
$x_i$ object of $\mathcal{X}$. Say $x_i$ and $y_i$ lie over the scheme $U_i$.
Then $f^*\mathcal{G}$ being quasi-coherent, means that $x_i^*f^*\mathcal{G}$
is quasi-coherent. Since $x_i^*f^*\mathcal{G}$ is isomorphic to
$y_i^*\mathcal{G}$ (as sheaves on $(\Sch/U_i)_{fppf}$ we
see that $y_i^*\mathcal{G}$ is quasi-coherent.
It follows from
Modules on Sites, Lemma \ref{sites-modules-lemma-local-final-object}
that the restriction of $\mathcal{G}$ to $\mathcal{Y}/y$ is
quasi-coherent. Hence $\mathcal{G}$ is quasi-coherent by
Lemma \ref{lemma-characterize-quasi-coherent}.
\end{proof}

\end{document}
